% ============================================================
%  MASLD & Stroke Risk — Meta-Analysis Manuscript
%  Title: Metabolic dysfunction–associated steatotic liver disease
%         and stroke risk: A systematic review and meta-analysis
%         exploring AF-associated and fibrosis-related pathways
% ============================================================
\documentclass[12pt,a4paper]{article}

% ---- Packages ----
\usepackage[utf8]{inputenc}
\usepackage[T1]{fontenc}
\usepackage{amsmath}
\usepackage{booktabs}
\usepackage{caption}
\usepackage{subcaption}
\usepackage{graphicx}
\usepackage{geometry}
\geometry{margin=1in}
\usepackage{setspace}
\doublespacing
\usepackage{lineno}
\linenumbers
\usepackage{fancyhdr}
\usepackage{hyperref}
\hypersetup{
  colorlinks=true,
  linkcolor=blue,
  citecolor=blue,
  urlcolor=blue
}
\usepackage{natbib}
\bibliographystyle{unsrtnat}
\usepackage{xcolor}
% \usepackage{multirow}      % Not in BasicTeX — install if needed
% \usepackage{enumitem}     % Not in BasicTeX — install if needed

% ---- Metadata ----
\title{\textbf{Metabolic dysfunction--associated steatotic liver disease and stroke risk:\\
A systematic review and meta-analysis exploring AF-associated and fibrosis-related pathways}}

\author{
  % Authors to be filled
}

\date{July 1, 2026}

\begin{document}
\maketitle

% ============================================================
%  ABSTRACT
% ============================================================
\begin{center}
\textbf{Abstract}
\end{center}
\vspace{-8pt}

\textbf{Background:} Metabolic dysfunction--associated steatotic liver disease (MASLD) has been linked to increased cardiovascular risk, but its specific association with stroke — and the potential pathways underlying this association — remains incompletely characterized. We conducted a systematic review and meta-analysis to quantify the MASLD--stroke association and explore two candidate pathways: atrial fibrillation (AF) and hepatic fibrosis.

\textbf{Methods:} We searched [databases] for studies reporting stroke risk in MASLD versus non-MASLD populations. Random-effects meta-analysis was performed (REML estimator). Heterogeneity was explored via meta-regression (AF adjustment as moderator), subgroup analyses, and fibrosis-stratified estimates. Evidence for an AF-associated pathway was evaluated across five hierarchical levels (Table S1). Fibrosis-related risk was assessed through severity-stratified estimates and longitudinal fibrosis dynamics.

\textbf{Results:} Across 21 studies including approximately 29.8 million participants, MASLD was associated with a 37\% increased risk of stroke (HR 1.37, 95\% CI 1.28--1.46), with substantial heterogeneity (I\textsuperscript{2}=98.6\%). Meta-regression identified AF adjustment as a significant moderator of between-study heterogeneity (R\textsuperscript{2}=16.4\%, p=0.040), with predicted stroke risk attenuating from HR 1.40 (95\% CI 1.31--1.49) to HR 1.13 (95\% CI 0.93--1.37) after AF adjustment. Patient-level interaction analysis (Jang 2026) confirmed effect modification by baseline AF (without AF: HR 1.10; with AF: HR 1.03, 0.95--1.12). Among four studies reporting fibrosis-stratified estimates, higher fibrosis burden was consistently associated with greater stroke risk, with the strongest signal from NFS-based stratification (Chen 2023: advanced fibrosis HR 2.08). Fibrosis regression was associated with reduced stroke risk (Park 2022: HR 0.58) compared to persistent fibrosis (HR 1.31).

\textbf{Conclusions:} MASLD is associated with increased stroke risk. These findings support the existence of an AF-associated pathway linking MASLD and stroke risk; however, formal mediation cannot be established from the available evidence. Fibrosis severity and dynamics may provide prognostic information beyond AF. These observations highlight the potential value of AF surveillance and fibrosis assessment in MASLD patients for stroke risk stratification.

\newpage

% ============================================================
%  INTRODUCTION
% ============================================================
\section{Introduction}

Metabolic dysfunction--associated steatotic liver disease (MASLD), formerly known as non-alcoholic fatty liver disease (NAFLD), affects approximately 25--30\% of the global adult population and represents the hepatic manifestation of the metabolic syndrome \cite{}. While MASLD is strongly associated with cardiovascular morbidity and mortality, its specific relationship with cerebrovascular events — particularly stroke — has received less systematic attention.

Several mechanisms have been proposed to link MASLD with stroke risk. Atrial fibrillation (AF), which is more prevalent in MASLD patients, may represent one such pathway: MASLD-associated systemic inflammation, autonomic dysfunction, and atrial structural remodeling may promote AF, which in turn increases thromboembolic stroke risk. Hepatic fibrosis severity may independently contribute to a prothrombotic state through altered coagulation factor synthesis and endothelial dysfunction.

However, existing primary studies vary substantially in their adjustment for AF and assessment of fibrosis, contributing to marked between-study heterogeneity. We conducted this systematic review and meta-analysis to: (1) quantify the overall association between MASLD and stroke risk; (2) evaluate the evidence for an AF-associated pathway linking MASLD and stroke; and (3) characterize the relationship between hepatic fibrosis severity/dynamics and stroke risk.

% ============================================================
%  METHODS
% ============================================================
\section{Methods}

\subsection{Search Strategy and Selection Criteria}

[Search strategy to be detailed.]

\subsection{Data Extraction and Quality Assessment}

[Data extraction procedures to be detailed.]

\subsection{Statistical Analysis}

The primary analysis used random-effects meta-analysis with restricted maximum-likelihood (REML) estimation to pool study-level hazard ratios. Heterogeneity was quantified using I\textsuperscript{2} and Cochran's Q. Meta-regression was performed to evaluate whether AF adjustment moderated the MASLD--stroke association. Subgroup analyses were conducted by fibrosis assessment status, AF adjustment, liver disease definition, and other study-level characteristics.

\textit{NHIS overlap handling.} Park et al. (2022) and Jang et al. (2026) originated from partially overlapping Korean National Health Insurance Service (NHIS) populations. To avoid double-counting, only Jang et al. (2026) was included in quantitative synthesis, whereas Park et al. (2022) was retained for descriptive evaluation of fibrosis dynamics.

All analyses were conducted in R (version 4.5.2) using the \texttt{meta} and \texttt{metafor} packages. Statistical significance was set at $\alpha = 0.05$ (two-sided).

% ============================================================
%  RESULTS
% ============================================================
\section{Results}

\subsection{Study Characteristics and Main Analysis}

The systematic search identified 2,192 records, of which 1,615 were screened after deduplication, 590 underwent full-text review, and 21 studies met inclusion criteria for quantitative synthesis (PRISMA flowchart, Supplementary Figure S1). Study characteristics are summarized in Table 1. The 21 included studies encompassed approximately 29.8 million participants from 5 countries (predominantly South Korea, n=14), published between 2019 and 2025. Study designs included population-based nationwide cohorts (n=8) and hospital- or clinic-based cohorts (n=13). MASLD was defined using NAFLD (n=8), MAFLD (n=7), or MASLD (n=6) criteria, with FLI-based ultrasonography as the most common diagnostic method. Median follow-up ranged from 3.2 to 15.8 years. Two studies adjusted for atrial fibrillation (AF); seven studies assessed hepatic fibrosis (Supplementary Table S2).

Across 21 studies, MASLD was associated with a 37\% increased risk of stroke (random-effects HR 1.37, 95\% CI 1.28--1.46, p$<$0.001; common-effect HR 1.32, 95\% CI 1.31--1.33; Figure 2). Individual study hazard ratios ranged from 1.10 (Jang 2026) to 2.01 (Kim 2020), all exceeding 1.0. Substantial between-study heterogeneity was observed (I\textsuperscript{2}=98.6\%, $\tau^2$=0.022, Cochran's Q=1403.9, df=20, p$<$0.001).

\subsection{Subgroup Analyses}

Subgroup analyses by stroke subtype showed comparable risk estimates for total stroke (k=13, HR 1.34, 95\% CI 1.26--1.43) and ischemic stroke (k=8, HR 1.40, 95\% CI 1.21--1.60), with no significant between-group difference (p=0.62; Supplementary Figure S2).

By liver disease definition, pooled estimates were lower for MASLD-defined studies (k=6, HR 1.20, 95\% CI 1.12--1.29) than for NAFLD-defined (k=8, HR 1.42, 95\% CI 1.23--1.63) or MAFLD-defined studies (k=7, HR 1.41, 95\% CI 1.34--1.48; p-between=0.002; Supplementary Figure S7).

By AF adjustment status, studies without AF adjustment (k=19) yielded a pooled HR of 1.40 (95\% CI 1.30--1.50), whereas the two studies adjusting for AF yielded a substantially lower pooled estimate (k=2, HR 1.11, 95\% CI 1.07--1.16; p-between$<$0.001; Supplementary Figure S8). However, this subgroup comparison is limited by the small number of AF-adjusted studies.

By setting, population-based studies (k=8, HR 1.42, 95\% CI 1.23--1.64) and hospital/clinic-based studies (k=13, HR 1.33, 95\% CI 1.26--1.40) yielded comparable estimates (p-between=0.40; Supplementary Figure S10).

\subsection{Meta-Regression}

Of seven study-level covariates evaluated, only AF adjustment was a statistically significant moderator of between-study heterogeneity in meta-regression (R\textsuperscript{2}=16.4\%, p=0.040; $\beta = -0.214$; Figure 6A). MASLD definition criteria showed a numerically large but non-significant contribution (R\textsuperscript{2}=13.8\%, p=0.139). The remaining covariates — diagnosis method (R\textsuperscript{2}=0.0\%, p=0.68), country (R\textsuperscript{2}=0.0\%, p=0.92), follow-up duration (R\textsuperscript{2}=0.0\%, p=0.91), sample size (R\textsuperscript{2}=1.8\%, p=0.28), and mean age (R\textsuperscript{2}=6.9\%, p=0.25) — did not explain significant proportions of between-study variance (Supplementary Figures S4--S6).

\subsection{AF-Associated Pathway}

Evidence for an AF-associated pathway was evaluated using a five-level hierarchy (Supplementary Table S1).

\textit{Meta-regression (Level I).} Meta-regression identified AF adjustment as a significant contributor to between-study heterogeneity (R\textsuperscript{2}=16.4\%, p=0.040; $\beta = -0.214$; Figure 6A). The predicted stroke risk, estimated from the regression model, was HR 1.40 (95\% CI 1.31--1.49) without AF adjustment compared to HR 1.13 (95\% CI 0.93--1.37) with AF adjustment, corresponding to 67.7\% attenuation of the excess risk attributable to AF adjustment.

\textit{Interaction by baseline AF (Level II).} In the only study reporting AF-stratified individual-level estimates (Jang 2026), the association between MASLD and stroke remained significant among individuals without AF (HR 1.10, 95\% CI 1.08--1.11) but was attenuated among those with pre-existing AF (HR 1.03, 95\% CI 0.95--1.12; Figure 6B), consistent with effect modification by baseline AF status. Within-study AF covariate adjustment yielded consistent directional attenuation (HR 1.14 $\rightarrow$ 1.10; 28.6\% attenuation).

\textit{Supporting pathway evidence (Level S1).} Ohno et al. (2023) demonstrated that MASLD was associated with a 51\% increased risk of incident AF (HR 1.51, 95\% CI 1.46--1.57; n=8.2 million), establishing the MASLD--AF link in the proposed pathway.

\textit{Additional evidence.} Lee et al. (2021) reported attenuation after bundle adjustment for cardiovascular comorbidities including AF (HR 1.55 $\rightarrow$ 1.43; 21.8\%), though the AF-specific contribution could not be isolated.

\subsection{Fibrosis-Related Pathway}

Four studies reported fibrosis-stratified stroke risk estimates using different assessment tools (histology, NFS, AAR, BARD). Individual study estimates are shown in Figure 7A. Across studies using histologic and non-invasive fibrosis assessment tools, higher fibrosis burden was consistently associated with greater stroke risk. The highest risk estimates were generally observed in advanced fibrosis categories, particularly when assessed by histology or NFS: histology (Simon 2022: low HR 1.50, intermediate HR 1.96, advanced HR 1.91), NFS (Chen 2023: low HR 1.29, intermediate HR 1.38, advanced HR 2.08), AAR (Chung 2022: reference HR 1.00, intermediate HR 1.13, high HR 1.29), and BARD (Jang 2026: low/reference HR 1.03, advanced HR 1.11).

Longitudinal fibrosis dynamics from Park et al. (2022) provided complementary evidence: fibrosis regression was associated with substantially reduced stroke risk (HR 0.58) compared to persistent fibrosis (HR 1.31; Figure 7B).

Although subgroup differences according to fibrosis assessment did not reach conventional statistical significance (p=0.078), studies evaluating more advanced fibrosis phenotypes tended to report higher risk estimates (Assessed: k=7, HR 1.26, 95\% CI 1.15--1.39; Not Assessed: k=14, HR 1.42, 95\% CI 1.30--1.55; Supplementary Figure S9). Exploratory pooled estimates by fibrosis severity category are presented in Supplementary Material only, as fibrosis assessment tools were not harmonized across studies.

\subsection{Sensitivity Analyses and Publication Bias}

The association remained robust across sensitivity analyses. Excluding Kim 2020 — the study with the largest effect estimate — yielded a pooled HR of 1.34 (95\% CI 1.26--1.41; k=20). Excluding Lee 2025 and Kim 2024, which shared NHIS source populations, yielded HR 1.37 (95\% CI 1.27--1.47; k=19). Leave-one-out analysis confirmed that no single study drove the overall estimate.

Publication bias assessment was inconclusive. Egger's test was non-significant in the full sample (p=0.64) but became significant after excluding the largest study (Egger p=0.018 without Lee 2025). Begg's rank correlation test was non-significant (p=1.00). Funnel plot asymmetry emerged after removal of Lee 2025, warranting cautious interpretation (Supplementary Figure S3).

% ============================================================
%  DISCUSSION
% ============================================================
\section{Discussion}

This systematic review and meta-analysis of 21 studies including approximately 29.8 million participants demonstrates that MASLD is associated with a 37\% increased risk of stroke. Beyond confirming this association, our analysis identifies two factors that contribute to risk heterogeneity: an AF-associated pathway and fibrosis severity/dynamics.

Meta-regression identified adjustment for atrial fibrillation (AF) as a significant source of between-study heterogeneity, explaining 16.4\% of the observed variance (R\textsuperscript{2}=16.4\%, p=0.040). Studies adjusting for AF generally reported lower effect estimates than those without AF adjustment, with predicted stroke risk attenuating from HR 1.40 to 1.13 after AF adjustment. This between-study finding is corroborated by patient-level interaction data (Jang 2026: HR 1.10 without baseline AF vs 1.03 with AF) and supported by large-scale evidence linking MASLD to incident AF (Ohno 2023: HR 1.51, n=8.2M). Taken together, these findings support the existence of an AF-associated pathway linking MASLD and stroke risk. However, formal mediation cannot be established from the available evidence, which derives from observational data and ecological meta-regression rather than individual-level mediation analysis. We therefore describe this as an ``AF-associated pathway'' rather than confirming AF mediation.

Regarding fibrosis, across studies using histologic and non-invasive fibrosis assessment tools, higher fibrosis burden was consistently associated with greater stroke risk. The highest risk estimates were observed in advanced fibrosis categories, particularly when assessed by histology or NFS (e.g., Chen 2023: advanced fibrosis HR 2.08). The dynamic nature of fibrosis-related risk — fibrosis regression associated with substantially reduced stroke risk (HR 0.58) compared to persistent fibrosis (HR 1.31) — further suggests that hepatic fibrosis may represent a modifiable risk modifier rather than a static confounder. Most included studies used FLI-based definitions of steatosis rather than fibrosis assessment tools; consequently, fibrosis-specific analyses were restricted to studies employing histology or validated fibrosis scores such as NFS and BARD.

Importantly, fibrosis-associated stroke risk remained evident even after adjustment for AF in Jang et al. (2026), suggesting that fibrosis-related cerebrovascular risk may not be fully explained by AF-associated mechanisms. These findings support a model in which AF-related and fibrosis-related pathways contribute independently to stroke risk in MASLD. However, although subgroup differences according to fibrosis assessment showed a consistent direction (assessed: HR 1.26; not assessed: HR 1.42), they did not reach conventional statistical significance (p=0.078), and the heterogeneity in fibrosis measurement methods (histology, NFS, AAR, BARD) precludes formal meta-analytic pooling.

Several limitations warrant consideration. First, the between-study meta-regression (Level I evidence) is ecological by design and cannot establish individual-level mediation. Second, only two studies in our sample adjusted for AF, limiting the statistical power of subgroup comparisons. Third, fibrosis-stratified estimates derive from single studies using non-standardized tools. Fourth, residual confounding by metabolic comorbidities shared between MASLD and stroke cannot be excluded.

In conclusion, MASLD is associated with increased stroke risk, with AF-associated mechanisms and fibrosis severity/dynamics representing potential contributors to risk heterogeneity. These findings support AF surveillance and fibrosis assessment as components of stroke risk stratification in MASLD patients.

% ============================================================
%  REFERENCES
% ============================================================
\newpage
\begin{thebibliography}{99}

% References will be added via BibTeX or manually

\end{thebibliography}

% ============================================================
%  FIGURE LEGENDS
% ============================================================
\newpage
\section*{Figure Legends}

\textbf{Figure 2.} Forest plot of the association between MASLD and stroke risk (main analysis). Random-effects meta-analysis (REML estimator), k=21, n$\approx$29.8M. Squares represent study-specific hazard ratios; size is proportional to study weight. Diamond represents the pooled estimate (HR 1.37, 95\% CI 1.28--1.46). I\textsuperscript{2}=98.6\%.

\textbf{Figure 6.} Evidence for an AF-associated pathway linking MASLD and stroke risk. \textbf{(A)} Meta-regression evidence: predicted hazard ratios from random-effects meta-regression with AF adjustment as moderator (R\textsuperscript{2}=16.4\%, p=0.040). Diamond vertices indicate predicted HRs and 95\% CIs. \textbf{(B)} Patient-level AF subgroup interaction (Jang 2026): dual forest plot without causal arrows. Effect modification annotation replaces directional arrow notation. \textbf{(C)} Left: mechanism cascade (MASLD $\rightarrow$ incident AF $\rightarrow$ stroke risk). Right: evidence chain boxes (meta-regression, AF subgroup interaction, MASLD$\rightarrow$AF pathway link).

\textbf{Figure 7.} Fibrosis severity and dynamics. \textbf{(A)} Fibrosis-stratified stroke risk estimates across assessment tools (histology, NFS, AAR, BARD). Shape coding: circle = low/reference, square = intermediate, triangle = advanced/high. Color coding by assessment tool. \textbf{(B)} Fibrosis dynamics (Park 2022): fibrosis regression (HR 0.58, protective) vs persistent fibrosis (HR 1.31).

\textbf{Supplementary Table S1.} Evidence hierarchy for the AF-associated pathway between MASLD and stroke risk. Five evidence levels: I (between-study meta-regression), II (patient-level interaction), III (within-study AF adjustment), IV (bundle adjustment), S1 (supporting pathway evidence). Attenuation calculated as (HR\textsubscript{without AF} $-$ HR\textsubscript{with AF}) / (HR\textsubscript{without AF} $-$ 1) $\times$ 100\%.

\textbf{Supplementary Figure S9.} Subgroup analysis by fibrosis assessment status (assessed vs not assessed). Random-effects model. p-between = 0.078.

\end{document}
